% !TeX spellcheck = en_us
\documentclass[a4paper,12pt,fleqn]{article}

\usepackage[left=1.50cm, right=1.50cm, top=3cm, bottom=2cm, 
headheight= 2.5cm, headsep= 4pt, footskip=25pt]{geometry}
\usepackage{fancyhdr}
\usepackage{lmodern}
\usepackage[table]{xcolor}
\usepackage{amssymb,amstext,amsmath}
\usepackage{booktabs}
\usepackage{bm}
\usepackage{setspace}
\usepackage{graphicx}
\usepackage[colorlinks]{hyperref}
\setlength{\parindent}{0mm}
\setlength{\parskip}{2mm}
\usepackage{enumitem}
\usepackage{titlesec}
\usepackage{gensymb}
\usepackage{pgfgantt}
\usepackage{lscape}
\usepackage[absolute]{textpos}
\usepackage{comment}
\usepackage{cite}
\usepackage{tikz}
\usepackage{pgfplots}
\usetikzlibrary{shapes.geometric}
\usepackage{array}
\newcolumntype{P}[1]{>{\centering\arraybackslash}p{#1}}
\usepackage{longtable}
\usepackage[font={small}]{caption}
\usepackage[most]{tcolorbox}
\usepackage{tikz}
\usepackage{stmaryrd}
\usepackage{wrapfig}
%\usepackage[none]{hyphenat}
\usepackage{ctable}
\usepackage{ragged2e}
\usepackage[official]{eurosym}
\newcommand*\circled[1]{\tikz[baseline=(char.base)]{
		\node[shape=circle,draw,inner sep=1pt] (char) {#1};}}
\usepackage[symbol]{footmisc}
\renewcommand{\thefootnote}{\fnsymbol{footnote}}
\renewcommand{\thempfootnote}{\fnsymbol{mpfootnote}}
\renewcommand*\footnoterule{}

\usepackage{xcolor}
\definecolor{docColor}{cmyk}{1.00,1.00,0, 0.40}  % 0.90\% of black
\color{docColor}

\newtcolorbox{myhlbox}[2][]{colback=docColor!05, colframe=docColor!20, width=1.0\linewidth,
	boxrule=0pt, left=1em, right=1em, top=1em, bottom=1em,
	colbacktitle=docColor!20, coltitle=docColor,enhanced, attach boxed title to top left={xshift=2mm, yshift=-1mm},
	title=#2, #1}

\hypersetup{
	colorlinks=true,
	citecolor=blue,
	filecolor=magenta,
	linkcolor=blue,
	urlcolor=blue,
	pdftitle={SeminarAnnouncements},
	%bookmarks=true,,
}

\graphicspath{{./figs/}}

\definecolor{dgreen}{rgb}{0.0, 0.5, 0.0}

\allowdisplaybreaks

\newcommand*\mystrut[1]{\vrule width0pt height0pt depth#1\relax}
\renewcommand{\thesection}{\arabic{section}.}
\renewcommand{\thesubsection}{\arabic{section}.\arabic{subsection}}
\renewcommand{\thesubsubsection}{\arabic{section}.\arabic{subsection}.\arabic{subsubsection}}
\titleformat{\section}{\sffamily\large\bfseries}{\thesection}{1.0em}{}
\titleformat{\subsection}{\sffamily\large\bfseries}{\thesubsection}{1.0em}{}
\titleformat{\subsubsection}{\sffamily\normalsize\bfseries}{\thesubsubsection}{1.0em}{}
\titlespacing\section{0pt}{0pt plus 2pt minus 2pt}{0pt plus 2pt minus 2pt}
\renewcommand{\baselinestretch}{1.20}

\newcommand\CircledImage[1]{%
	\tikz\path[fill overzoom image={#1}, line width=0mm]circle[radius=0.5\linewidth];%
}

\input{format/defs}

\fancypagestyle{titlepage}{
	\fancyhf{}
	\insertLXheader
}

\fancypagestyle{plain}{%
	\fancyhf{}
	\renewcommand{\headrulewidth}{0pt}
	\renewcommand{\footrulewidth}{0pt}
	\insertnormalpage
}

\pagestyle{plain}

\begin{document}
	\sffamily
\thispagestyle{titlepage}
\vspace*{-3em}
\begin{center}
	\Large \textbf{LMS internal discussion}
\end{center}
\begin{center}
	\Large 
	\textbf{Fabrication across scales for mechanics and materials research:~Current capabilities and future needs} 
\end{center}
\begin{center}
	\small
	\vspace*{-0.5em}
	\textbf{Moderator:}~Vignesh Kannan;	\textbf{Speakers:}~Alistair Rowe (LPMC - Laboratoire de Physique de la Matière Condensée), Andrea Rocchi (Drahi'X Innovation center), Laurence Bodelot \& Vignesh Kannan (LMS) \\
	\vspace*{0.5em}
	{\large
	\textbf{Date and Time}:~September 17, 2026 (2 -- 3 pm)
	
	\textbf{Venue}:~Amphi 104 (P{\^o}le M{\'e}ca) }
\end{center}
\vspace*{-1em}
\begin{myhlbox}{{\textbf{Abstract}}}
	This internal seminar will focus on the current capabilities (and needs) at Polytechnique, for fabrication across scales from materials and microstructure, to precise and complex sample geometries.

	In the realm of \textbf{sample fabrication}, conventional machining has proved to be reliable above the millimeter length scales. Complexities in sample geometry may be achieved using multi-axis milling, with recent progress allowing us to access sub-millimeter scales. However, fabrication at this intermediate scale still remains a challenge in our labs. At the micro- and nano-scales, photo- and e-beam lithography, silicon and metal etching techniques could prove powerful to push the limits of sample design, taking us to the realm of \textbf{materials fabrication}. These microfabrication techniques are scalable with the potential for excellent microstructural control, potentially allowing multiscale control of design parameters from the material to sample/structure. Polymer-based 3D printing and lithography have also made significant progress in parallel, to enable multifunctionality and small-scale testing of soft and biological media. Taking these capabilities even further down in scales, chemical composition control for thin films is an achievable goal in cleanrooms these days, and additive manufacturing is attempting to control alloy compositions across larger volumes.

	The convergence of modern techniques to make ``materials'' and ``samples'' is especially crucial given modern directions in the experimental sciences, and the interplay of experiment, theory, computation and design for microstructural and multifunctional control of materials across scales. 
\end{myhlbox}

\begin{myhlbox}{\textbf{Speakers and Format}}
	This discussion will start with \textbf{Prof.~Alistair Rowe}, professor at l'X and researcher at LPCM, who is also leading the development of a new cleanroon microfabrication facility at l'X (MicroFabriX). His talk will be followed by \textbf{Prof.~Laurence Bodelot} (and Vignesh) on fabrication capabilities within the LMS, followed by \textbf{Mr. Andrea Rocchi}, technical manager at the Drahi Innovation center of l'X on larger scale manufacturing infrastructure. 

	Each speaker will have 10~mins for presentation followed by a 3-5~min Q\&A. The last 10 - 15~minutes will be dedicated to a group discussion about the using these capabilities for new research directions, and identifying gaps in fabrication length scales to be addressed.  
\end{myhlbox}

\end{document}